\section{Bayesian Point Calibration and Zero-Inflated DNF Penalties}
\label{sec:bayesian_point_calibration}

In fantasy sports, winning portfolios are determined not by raw finishing positions, but by the expected value and variance of rule-based fantasy points $\text{Pts} \in \mathbb{R}$. Formula 1 Fantasy employs a highly non-linear, multi-attribute scoring schedule that penalizes retirements severely.

\subsection{Official Scoring Function and Asymmetric Payoffs}
For driver $i$ with starting grid $g_i$ and finishing position $p_i$, the deterministic fantasy point reward is:
\begin{align}
\text{Pts}(g_i, p_i) = &\; \text{Pts}_{\text{race}}(p_i) + \text{Pts}_{\text{quali}}(g_i) \nonumber \\
&+ \text{clamp}\left( g_i - p_i, -10, +10 \right) \nonumber \\
&+ \text{Bonus}_{\text{teammate}} + \text{Bonus}_{\text{FL}} + \text{Bonus}_{\text{DOTD}}
\label{eq:scoring_formula}
\end{align}
where $\text{Pts}_{\text{race}}(p_i)$ follows the FIA World Championship scale ($25, 18, 15, \dots$), positions gained/lost award $\pm 1$~point per position delta (bounded within $[-10, +10]$), and defeating a teammate awards $+3$~points.

Critically, a Did Not Finish (DNF) event incurs a direct penalty:
\begin{equation}
\text{Pts}_{\text{DNF}} = -15.0 \text{ points}
\label{eq:dnf_penalty}
\end{equation}

\subsection{The Zero-Inflation Retirement Bias and Fix}
Prior to Phase 2 of our verification audit, standard prediction engines suffered from \textit{zero-inflation optimistic bias}. Specifically, when computing expected points $\mathbb{E}[\text{Pts}_i]$, models either omitted DNF outcomes or mapped retirements to $p_i = 20$ (which yielded $0$ points plus position-loss penalties). Consequently, high-risk backmarker drivers with $25\%$ mechanical retirement probability were assigned positive expected yields ($\mathbb{E} \approx +6.5$ pts), triggering disastrous algorithmic allocations in low-budget optimizer rosters.

To establish rigorous calibration, we formulated a zero-inflated two-component mixture distribution \cite{gelman2013bayesian}. Let $Z_i \in \{0, 1\}$ denote the binary retirement indicator ($Z_i = 1$ denotes DNF) with probability $P_{\text{dnf}, i}$. The true calibrated expected fantasy point yield is:
\begin{equation}
\mathbb{E}[\text{Pts}_i] = (1 - P_{\text{dnf}, i}) \cdot \mathbb{E}[\text{Pts}_i \mid Z_i = 0] + P_{\text{dnf}, i} \cdot \text{DNF\_PENALTY}
\label{eq:calibrated_expectation}
\end{equation}

The retirement probability $P_{\text{dnf}, i}$ is modeled via empirical Bayes pooling across driver, chassis, and circuit risk factors:
\begin{equation}
P_{\text{dnf}, i} = \text{clamp}\left( 0.40 \cdot \theta_{\text{driver}, i} + 0.40 \cdot \theta_{\text{team}, c} + 0.20 \cdot \theta_{\text{circuit}}, 0.02, 0.45 \right)
\label{eq:dnf_probability}
\end{equation}
where $\theta_{\text{driver}, i}$ and $\theta_{\text{team}, c}$ represent historical 2-season rolling DNF rates, and $\theta_{\text{circuit}}$ represents track incident frequency (e.g., Singapore, Baku, or Monaco).

Under Eq.~\eqref{eq:calibrated_expectation}, volatile drivers with elevated retirement rates appropriately register negative or depressed expected yields ($\mathbb{E}[\text{Pts}] \approx -2.0$ to $+1.5$), protecting the downstream optimizer from fragile assets.

\subsection{Conjugate Beta-Binomial Teammate Dominance}
The teammate bonus ($+3$ pts) represents an important source of alpha. For teammates $(A, B)$, the probability $P(A > B)$ is modeled as a Beta-Binomial conjugate posterior:
\begin{equation}
P(A > B \mid \mathcal{D}) = \frac{\alpha_0 + k_A}{\alpha_0 + \beta_0 + N_{AB}}
\label{eq:beta_binomial}
\end{equation}
where $\alpha_0 = \beta_0 = 2$ represents an uninformative prior, $N_{AB}$ denotes completed head-to-head races, and $k_A$ denotes victories by driver $A$.
